\chapter{Complexity — Step}
%%% generated by the blueprint pipeline from TCSlib.Complexity.SpaceComplexity.Codec.Step
\section{Overview}

Declarations of the Lean module \texttt{TCSlib.Complexity.SpaceComplexity.Codec.Step}.

\section{Declarations}

\begin{definition}[Idle machine action]
\lean{Complexity.Codec.idleAct}
\label{Complexity.Codec.idleAct}
\leanok
\uses{Turing.Action}
The idle action leaves the input head and every work-tape head stationary, makes no
work-tape writes, emits no output symbol, and has no successor state.
\end{definition}

\begin{definition}[Effective machine action]
\lean{Complexity.Codec.effAct}
\label{Complexity.Codec.effAct}
\leanok
\uses{Complexity.Codec.idleAct, Turing.Action}
Given an optional current state, an optional input symbol, the symbols scanned on the
work tapes, and a choice bit, the effective action is the idle action when there is no
current state. When there is a current state, it is the machine's transition action for
that state, scanned symbols, and choice bit.
\end{definition}

\begin{theorem}[Choice Step as an Applied Action]
\lean{Complexity.Codec.stepWith_eq_apply}
\label{Complexity.Codec.stepWith_eq_apply}
\leanok
\uses{Complexity.Codec.effAct, Complexity.Codec.idleAct, Turing.Action.apply, Turing.Cfg, Turing.Cfg.inputSymbol, Turing.Cfg.workTapeSymbols, Turing.NDTM.stepWith, Turing.NDTM.stepWith_of_halt, Turing.moveInputPos_zero}
\difficulty{4}
Let $N$ be a finite nondeterministic Turing machine with Boolean symbols, let $x$ be an
input, let $b\in\{0,1\}$, and let $C$ be a configuration on $x$. The result of taking
one step from $C$ using choice $b$ is the configuration obtained by applying to $C$ the
action prescribed by the machine for choice $b$, using the current state, the symbol
under the input head, and the symbols under the work-tape heads; when $C$ is halted,
this prescribed action is the idle action.
\end{theorem}

\begin{definition}[Output summary update]
\lean{Complexity.Codec.sumUpd}
\label{Complexity.Codec.sumUpd}
\leanok
Define an update on the three summary states $0,1,2$ for an optional Boolean output.
From state $0$, no output leaves the state at $0$, outputting true changes it to $1$,
and outputting false changes it to $2$; from state $1$, no output leaves it at $1$ and
either Boolean output changes it to $2$; state $2$ remains $2$ for every output.
\end{definition}

\begin{lemma}[Output summary under appending a bit]
\lean{Complexity.Codec.sumIdx_append}
\label{Complexity.Codec.sumIdx_append}
\leanok
\uses{Complexity.Codec.sumIdx, Complexity.Codec.sumUpd, Turing.outSummary}
\difficulty{4}
For any Boolean list $u$ and optional Boolean symbol $o$, the summary of the output
obtained by appending $o$ to $u$ is the update of the summary of $u$ by $o$. Here the
three summaries distinguish an empty output, the output consisting exactly of the single
bit $1$, and every other output.
\end{lemma}

\begin{definition}[Updated head summary]
\lean{Complexity.Codec.newHd}
\label{Complexity.Codec.newHd}
\leanok
\uses{Complexity.Codec.Hd, Complexity.Codec.sumUpd, Turing.Action}
For a machine action and a head summary consisting of an optional state and a value in
$0,1,2$, the updated summary has the action's optional successor state and the result
of applying the output-summary update to the old value and the action's optional Boolean
output.
\end{definition}

\begin{definition}[Blank cell encoding]
\lean{Complexity.Codec.cellE}
\label{Complexity.Codec.cellE}
\leanok
\uses{Complexity.Codec.CellV}
The blank cell encoding has no head summary, and on each work tape it records no symbol
and the Boolean value false.
\end{definition}

\begin{definition}[Local cell update]
\lean{Complexity.Codec.stepCell}
\label{Complexity.Codec.stepCell}
\leanok
\uses{Complexity.Codec.CellV, Complexity.Codec.Hd, Turing.Action}
Given a machine action, a new head summary, and left, central, and right cell encodings,
the updated cell carries the new head summary exactly when the central cell carried some
head summary; otherwise it carries none. For each work tape, it retains the central
cell's symbol unless its Boolean marker is true, in which case it uses the action's
optional write value with the old symbol as default; its new Boolean value is the
central marker for a stationary move, the left neighbor's marker for a positive move,
and the right neighbor's marker for a negative move.
\end{definition}

\begin{definition}[Partial local update]
\lean{Complexity.Codec.stepG}
\label{Complexity.Codec.stepG}
\leanok
\uses{Complexity.Codec.CellV, Complexity.Codec.Hd, Complexity.Codec.stepCell, Turing.Action}
Given a machine action and a new head summary, define a function on three optional cell
encodings: when all three are present, return the present encoding obtained by the local
cell update; in every other case, return none.
\end{definition}

\begin{definition}[Head summary of a configuration]
\lean{Complexity.Codec.hdOf}
\label{Complexity.Codec.hdOf}
\leanok
\uses{Complexity.Codec.Hd, Complexity.Codec.sumIdx, Turing.Cfg, Turing.outSummary}
For a configuration of a machine on a fixed Boolean input, its head summary consists of
its optional state and the value obtained by applying the output-summary encoding to its
output string and then mapping that summary to one of three values.
\end{definition}

\begin{definition}[Left boundary cell]
\lean{Complexity.Codec.nbL}
\label{Complexity.Codec.nbL}
\leanok
\uses{Complexity.Codec.CellV, Complexity.Codec.cellE, Complexity.Codec.cellOf, Turing.Cfg}
For a natural-number parameter, a machine configuration on a fixed Boolean input, and a
natural-number cell index, define the left neighbor encoding to be the blank cell
encoding when the index is zero; otherwise it is the configuration's cell encoding at
index one less than the given index.
\end{definition}

\begin{definition}[Right boundary cell]
\lean{Complexity.Codec.nbR}
\label{Complexity.Codec.nbR}
\leanok
\uses{Complexity.Codec.CellV, Complexity.Codec.cellE, Complexity.Codec.cellOf, Turing.Cfg}
For a natural-number parameter, a machine configuration on a fixed Boolean input, and a
natural-number cell index, define the right neighbor encoding to be the configuration's
cell encoding at index one greater than the given index when that index is less than
twice the parameter; otherwise it is the blank cell encoding.
\end{definition}

\begin{lemma}[Integer value of a sign]
\lean{Complexity.Codec.signType_cast}
\label{Complexity.Codec.signType_cast}
\leanok
\difficulty{4}
For each sign $m\in\{-1,0,1\}$, its integer cast is $0$ when $m$ is zero, $1$ when $m$
is positive, and $-1$ when $m$ is negative.
\end{lemma}

\begin{theorem}[Cell encoding under one step]
\lean{Complexity.Codec.cellOf_apply}
\label{Complexity.Codec.cellOf_apply}
\leanok
\uses{Complexity.Codec.cellE, Complexity.Codec.cellOf, Complexity.Codec.hdOf, Complexity.Codec.nbL, Complexity.Codec.nbR, Complexity.Codec.newHd, Complexity.Codec.signType_cast, Complexity.Codec.stepCell, Complexity.Codec.sumIdx_append, Turing.Action, Turing.Action.apply, Turing.Cfg, Turing.Cfg.InWindow}
\difficulty{5}
Let $N$ be a finite nondeterministic Turing machine with Boolean symbols, let $x$ be a
finite Boolean string, and let $s$ be a natural number. Suppose configuration $C$ has
all work-tape heads within distance $s$ of the origin and blank work-tape cells outside
that window. For any one-step action $a$, any natural number $j<2s+1$, and the updated
head-and-output summary determined by $a$ and $C$, the encoded cell at position $j$ of
the configuration after applying $a$ equals the local update of the encoded cell at $j$
and its two neighboring cells, using an empty cell beyond either endpoint of the window.
\end{theorem}

\begin{theorem}[Boundary test for window preservation]
\lean{Complexity.Codec.boundary_iff}
\label{Complexity.Codec.boundary_iff}
\leanok
\uses{Complexity.Codec.Hd, Complexity.Codec.cellE, Complexity.Codec.cellOf, Complexity.Codec.signType_cast, Complexity.Codec.stepCell, Turing.Action, Turing.Action.apply, Turing.Cfg, Turing.Cfg.InWindow}
\difficulty{5}
Let $N$ be a finite nondeterministic Turing machine with Boolean symbols, let $x$ be a
finite Boolean string, and let $s$ be a natural number. Suppose a configuration $C$ has
every work-tape head within distance $s$ of the origin and has blank cells at every
work-tape position whose distance from the origin exceeds $s$. For any one-step action
$a$ of the machine and any pair $h'$ consisting of an optional state and a three-valued
output summary, the following are equivalent: updating the encoded cell at the right
endpoint of the window using the empty cell on its left and right gives the empty cell,
and updating the encoded cell at the left endpoint using the empty cell on its left and
right gives the empty cell; the resulting configuration after action $a$ again has all
work-tape heads within distance $s$ and blank cells at every position farther than $s$
from the origin.
\end{theorem}

\begin{theorem}[Local Characterization of a Machine Step]
\lean{Complexity.Codec.slide_step_iff}
\label{Complexity.Codec.slide_step_iff}
\leanok
\uses{Complexity.Codec.CellV, Complexity.Codec.boundary_iff, Complexity.Codec.cellE, Complexity.Codec.cellOf, Complexity.Codec.cellOf_apply, Complexity.Codec.hdOf, Complexity.Codec.nbL, Complexity.Codec.nbR, Complexity.Codec.newHd, Complexity.Codec.stepG, Turing.Action, Turing.Action.apply, Turing.Cfg, Turing.Cfg.InWindow}
\difficulty{6}
Let $N$ be a finite nondeterministic Turing machine with Boolean input and $k$ work
tapes, let $x$ be an input, let $s\in\mathbb{N}$, and let $C$ be a configuration on $x$
whose work-tape heads lie within distance $s$ of the origin and whose tape contents
vanish outside that window. Let $a$ be an action and let $v$ be a list of optional cell
values of length $2s+1$. The condition that, at every position $i<2s+3$, the padded list
of values in $v$ agrees with the local update rule determined by $a$ and the three
corresponding padded cells of the encoding of $C$ is equivalent to both of the
following: each entry of $v$ is the encoded cell of the configuration obtained by
applying $a$ to $C$, and that resulting configuration has all work-tape heads and
nonblank contents within the same window of radius $s$.
\end{theorem}

\begin{definition}[Optional machine cell]
\lean{Complexity.Codec.SX}
\label{Complexity.Codec.SX}
\leanok
\uses{Complexity.Codec.CellV}
For a nondeterministic Turing machine $N$ with Boolean tape alphabet, let a machine cell
be a cell value of $N$ when present, and otherwise absent; thus the type of optional
machine cells is the option type of the cell values of $N$.
\end{definition}

\begin{definition}[Pair of scan states]
\lean{Complexity.Codec.SSt}
\label{Complexity.Codec.SSt}
\leanok
\uses{Complexity.Codec.SX, Complexity.Codec.SlSt}
For a finite nondeterministic Turing machine $N$ with Boolean tape alphabet, a paired
scan state consists of two scan states, each having three optional machine-cell entries
and a Boolean entry.
\end{definition}

\begin{definition}[Decoding paired chunks]
\lean{Complexity.Codec.pairChunk}
\label{Complexity.Codec.pairChunk}
\leanok
\uses{Complexity.Codec.SX, Complexity.Codec.decChunk}
Let $N$ be a finite nondeterministic Turing machine with Boolean tape alphabet, and let
$ch$ be a finite list of pairs of Boolean symbols. Decode the list of first coordinates
of $ch$ and the list of second coordinates of $ch$ separately into machine chunks; the
resulting pair of optional cell values is the decoded paired chunk.
\end{definition}

\begin{definition}[Initial scan state]
\lean{Complexity.Codec.sInit}
\label{Complexity.Codec.sInit}
\leanok
\uses{Complexity.Codec.SX, Complexity.Codec.SlSt, Complexity.Codec.cellE}
Let $N$ be a finite nondeterministic Turing machine with Boolean tape alphabet, and let
$e$ be its distinguished cell value. The initial scan state has all three cell entries
equal to the present value $e$ and its Boolean entry true.
\end{definition}

\begin{definition}[Paired scan transition]
\lean{Complexity.Codec.sChunk}
\label{Complexity.Codec.sChunk}
\leanok
\uses{Complexity.Codec.SSt, Complexity.Codec.SX, Complexity.Codec.pairChunk, Complexity.Codec.slStep}
Let $N$ be a finite nondeterministic Turing machine with Boolean tape alphabet, let
$G_0,G_1$ each take three optional cell values to an optional cell value, let $st$ be a
pair of scan states, and let $ch$ be a finite list of pairs of Boolean symbols. Decode
the two coordinate lists of $ch$ into a pair of optional cell values, then apply the
scan-state step determined by $G_0$ to the first scan state and by $G_1$ to the second,
using that decoded pair as input. The resulting pair of scan states is the paired scan
transition.
\end{definition}

\begin{definition}[Paired scan of two lists]
\lean{Complexity.Codec.sScan}
\label{Complexity.Codec.sScan}
\leanok
\uses{Complexity.Codec.SSt, Complexity.Codec.SX, Complexity.Codec.adStep, Complexity.Codec.hH, Complexity.Codec.sChunk, Complexity.Codec.sInit}
Let $N$ be a finite nondeterministic Turing machine with Boolean tape alphabet, let
$G_0,G_1$ each take three optional cell values to an optional cell value, and let $u,u'$
be finite lists of Boolean symbols. Pair the lists position by position, apply the
paired scan transition successively across the resulting list of pairs, starting from
two initial scan states together with counter $0$ and a function assigning
$(	ext{false},	ext{false})$ at every argument, and return the pair of scan states in the
resulting accumulator.
\end{definition}

\begin{definition}[Completion flag after two steps]
\lean{Complexity.Codec.sDone}
\label{Complexity.Codec.sDone}
\leanok
\uses{Complexity.Codec.SX, Complexity.Codec.SlSt, Complexity.Codec.cellE, Complexity.Codec.slStep}
Let $N$ be a finite nondeterministic Turing machine with Boolean tape alphabet, let $G$
take three optional cell values to an optional cell value, and let $st$ be a scan state.
Apply the scan-state step determined by $G$ twice, each time with the pair consisting of
two present copies of the distinguished cell value, and take the Boolean entry of the
resulting scan state. This Boolean is the completion flag after the two steps.
\end{definition}

\begin{definition}[Decoded blocks of a Boolean list]
\lean{Complexity.Codec.blocksD}
\label{Complexity.Codec.blocksD}
\leanok
\uses{Complexity.Codec.SX, Complexity.Codec.block, Complexity.Codec.decChunk}
Let $N$ be a finite nondeterministic Turing machine with Boolean tape alphabet, let $u$
be a finite list of Boolean symbols, and let $m$ be a natural number. For each integer
$j$ from $0$ through $m-1$, take the $j$th block of $u$ and decode it into an optional
machine-cell value; the ordered list of these decoded values is the list of decoded
blocks.
\end{definition}

\begin{lemma}[Folding paired updates componentwise]
\lean{Complexity.Codec.foldl_pair}
\label{Complexity.Codec.foldl_pair}
\leanok
\difficulty{4}
Let $\alpha,\beta,\gamma$ be types, and let $f : \alpha\to\gamma\to\alpha$ and $g :
\beta\to\gamma\to\beta$. For initial values $a\in\alpha$ and $b\in\beta$ and any finite
list $L$ over $\gamma$, folding the paired update on $(a,b)$ gives the pair consisting
of the separate folds of $f$ from $a$ and $g$ from $b$ over $L$.
\end{lemma}

\begin{theorem}[Blockwise decomposition of a paired scan]
\lean{Complexity.Codec.sScan_eq}
\label{Complexity.Codec.sScan_eq}
\leanok
\uses{Complexity.Codec.SX, Complexity.Codec.W, Complexity.Codec.block, Complexity.Codec.blocksD, Complexity.Codec.eq_flatten_blocks, Complexity.Codec.foldl_adStep_flatten, Complexity.Codec.foldl_pair, Complexity.Codec.hH, Complexity.Codec.pairChunk, Complexity.Codec.sChunk, Complexity.Codec.sInit, Complexity.Codec.sScan, Complexity.Codec.slStep, Complexity.Codec.zip_flatten}
\difficulty{6}
Let $N$ be a finite nondeterministic Turing machine over Boolean symbols, let $G_0,G_1$
be local update rules on optional encoded cells, and let $u,u'$ be Boolean words. If
both words have length $W_Nm$ for some natural number $m$, then scanning the paired
words one bit at a time and applying the two update rules produces the same pair of scan
states as pairing their decoded length-$W_N$ blocks and folding the respective local
update rules over those block pairs, starting from the initial scan state.
\end{theorem}

\begin{lemma}[Decoding blocks of a concatenated chunk word]
\lean{Complexity.Codec.blocksD_chunksWord}
\label{Complexity.Codec.blocksD_chunksWord}
\leanok
\uses{Complexity.Codec.CellV, Complexity.Codec.block_chunksWord, Complexity.Codec.blocksD, Complexity.Codec.chunksWord, Complexity.Codec.decChunk_encChunk}
\difficulty{3}
Let $N$ be a finite nondeterministic Turing machine over Boolean symbols, let $f$ assign
a cell value to each natural number, and let $m$ be a natural number. Decoding the first
$m$ blocks of the concatenation of the first $m$ encoded cell values gives, in order,
the values $f(0),\ldots,f(m-1)$, each represented as a successfully decoded cell value.
\end{lemma}

\begin{theorem}[Correctness of the step scan]
\lean{Complexity.Codec.sDone_step}
\label{Complexity.Codec.sDone_step}
\leanok
\uses{Complexity.Codec.W, Complexity.Codec.block_chunksWord, Complexity.Codec.blocksD, Complexity.Codec.blocksD_chunksWord, Complexity.Codec.cellE, Complexity.Codec.cellOf, Complexity.Codec.chunksWord, Complexity.Codec.encChunk_of_decChunk, Complexity.Codec.eq_flatten_blocks, Complexity.Codec.hdOf, Complexity.Codec.newHd, Complexity.Codec.sDone, Complexity.Codec.sInit, Complexity.Codec.slStep, Complexity.Codec.slide_ok_iff, Complexity.Codec.slide_step_iff, Complexity.Codec.stepG, Turing.Action, Turing.Action.apply, Turing.Cfg, Turing.Cfg.InWindow}
\difficulty{6}
Let $N$ be a finite nondeterministic Turing machine over Boolean symbols, let $x$ be a
Boolean word, let $s$ be a natural number, and let $C$ be a configuration of $N$ on
input $x$ whose work-tape heads and nonblank cells lie within distance $s$ of the
origin. Let $a$ be an action and let $A'$ be a Boolean word of length $W_N(2s+1)$. The
scan’s acceptance flag is true exactly when $A'$ is the encoded sequence of cells of the
configuration obtained by applying $a$ to $C$ across the window of width $2s+1$, and the
resulting configuration still has all work-tape heads and nonblank cells within distance
$s$ of the origin.
\end{theorem}
